\documentclass[11pt,a4paper]{article}

\usepackage[T1]{fontenc}
\usepackage[utf8]{inputenc}
\usepackage[brazilian]{babel}
\usepackage{lmodern}
\usepackage{microtype}
\usepackage[a4paper,margin=2.5cm]{geometry}
\usepackage{booktabs}
\usepackage{tabularx}
\usepackage{enumitem}
\usepackage{caption}
\usepackage{listings}
\usepackage{tikz}
\usetikzlibrary{automata,positioning,arrows.meta}
\usepackage[hidelinks]{hyperref}

% Números extraídos das evidências reais (gerado por gerar_relatorio.py)
% Gerado por relatorio/gerar_relatorio.py a partir de evidencias/. Não editar à mão.
\newcommand{\integrantes}{Alex Felippe Santana Guimarães \\ Gabriel Baeta Pereira \\ Lucas D'almeida Gavião \\ Vitor Lemos da Silva}
\newcommand{\numIndicesPublic}{22}
\newcommand{\numFKsPublic}{11}
\newcommand{\numQquatroTabelas}{11}
\newcommand{\numQquatroPKs}{11}
\newcommand{\numQquatroFKs}{11}
\newcommand{\numQquatroColunas}{65}
\newcommand{\numQquatroNotNull}{31}
\newcommand{\numCasosDTE}{55}
\newcommand{\numDTEa}{11}
\newcommand{\numDTEb}{11}
\newcommand{\numDTEc}{4}
\newcommand{\numDTEd}{23}
\newcommand{\numDTEe}{4}
\newcommand{\numDTEf}{2}
\newcommand{\numCasosH}{16}
\newcommand{\numCasosT}{24}
\newcommand{\numCasosS}{22}
\newcommand{\numCasosP}{7}
\newcommand{\numCasosSApp}{19}
\newcommand{\numVerifS}{3}
\newcommand{\numSessoes}{30}


\lstset{
  language=SQL,
  basicstyle=\ttfamily\small,
  keywordstyle=\bfseries,
  columns=fullflexible,
  keepspaces=true,
  upquote=true,
  showstringspaces=false,
  frame=lines,
  xleftmargin=1em,
  framexleftmargin=0.5em,
  aboveskip=0.8em,
  belowskip=0.8em,
  morekeywords={PROCEDURE,RETURNS,LANGUAGE,DEFINER,LATERAL,ORDINALITY,RECURSIVE,CALL,TRIGGER,EXECUTE,FUNCTION,FOR,EACH,ROW,BEFORE,AFTER,INOUT,OWNER,GRANT,REVOKE}
}

\setlist{itemsep=0.2em,topsep=0.4em}
% código em linha: permite quebra de linha após "_" e "." em identificadores longos
\newcommand{\cod}[1]{{\ttfamily\renewcommand{\_}{\textunderscore\allowbreak}#1}}
\emergencystretch=2em
\newenvironment{enunciado}{\begin{quote}\small\itshape}{\end{quote}}

\title{%
  \normalsize Universidade Federal Fluminense --- Instituto de Computação\\
  Banco de Dados II --- Prof.\ Rodrigo Salvador Monteiro\\[2.5em]
  \LARGE Trabalho em Grupo --- Partes 1 e 2\\[0.4em]
  \large Programação em banco de dados com PostgreSQL sobre o esquema Chinook}
\author{\integrantes}
\date{Niterói, \today}

\begin{document}

\maketitle
\thispagestyle{empty}

\begin{abstract}
\noindent Este relatório descreve as soluções do grupo para as Partes~1 e~2 do trabalho, implementadas em
PostgreSQL~16 com a linguagem procedural PL/pgSQL sobre o esquema exemplo Chinook. Para cada questão são
apresentados o enunciado, a solução adotada e sua justificativa, e o resultado obtido na execução. O código-fonte
completo, os testes automatizados e as saídas de execução acompanham este relatório no pacote de entrega; o texto
cita apenas os trechos necessários para explicar cada solução.
\end{abstract}

\tableofcontents

% =====================================================================
\section{Introdução}

\subsection{Ambiente}

O SGBD utilizado é o PostgreSQL~16, escolha aprovada pelo professor. O esquema exemplo é o Chinook Database
v1.4.5, na variante \cod{Chinook\_PostgreSql\_SerialPKs.sql}, que cria o banco \cod{chinook\_serial} com as onze
tabelas no schema \cod{public} e nomes em minúsculas separados por sublinhado (\cod{invoice}, \cod{invoice\_line},
\cod{unit\_price}, \dots). Essa variante foi escolhida porque suas chaves primárias têm valor padrão gerado por
sequência, sincronizada com os dados originais; assim, faturas e itens podem ser inseridos sem informar o
identificador, o que o procedimento da Parte~2 e os testes exigem.

As rotinas genéricas da Parte~1 e a configuração do motor de estados ficam no schema auxiliar \cod{bd2\_meta},
e os objetos de demonstração em \cod{bd2\_demo}. As regras da Parte~2, específicas do Chinook, ficam em
\cod{public}, junto às tabelas que protegem. Duas extensões do esquema foram feitas exclusivamente para demonstrar a Questão~5 da Parte~1: a
coluna \cod{invoice.status} e a tabela \cod{bd2\_demo.pedido\_demo}.

\subsection{Organização da entrega e execução}

\begin{table}[!htb]
\centering\small
\begin{tabularx}{\textwidth}{@{}lX@{}}
\toprule
Pasta / arquivo & Conteúdo \\
\midrule
\cod{sql/00\_ambiente/} & Script oficial do Chinook e verificação do ambiente \\
\cod{sql/01\_catalogos/} & Parte~1, Questões~1 e~3 \\
\cod{sql/02\_administracao/} & Parte~1, Questão~2 \\
\cod{sql/03\_ddl\_dinamico/} & Parte~1, Questão~4 \\
\cod{sql/04\_maquina\_estados/} & Parte~1, Questão~5 (motor genérico e cenário de demonstração) \\
\cod{sql/05\_regras\_semanticas/} & Parte~2 (papéis, triggers e procedimento armazenado) \\
\cod{sql/06\_testes/} & Testes automatizados, inclusive de concorrência \\
\cod{sql/07\_evidencias/} & Geração das saídas de cada questão \\
\cod{evidencias/} & Saídas da execução que embasa este relatório \\
\cod{run\_all.sh}, \cod{run\_all.ps1} & Execução completa a partir de um banco limpo \\
\bottomrule
\end{tabularx}
\caption{Organização do pacote de entrega.}
\end{table}

O script \cod{run\_all.sh} (ou \cod{run\_all.ps1}, no Windows) recria o banco a partir do script oficial do
Chinook, instala os módulos na ordem de dependência, executa todos os testes e grava as saídas em
\cod{evidencias/}. Os testes são executados pelo \cod{psql} com \cod{ON\_ERROR\_STOP}: qualquer divergência entre
resultado esperado e obtido interrompe a execução. Eles não dependem de identificadores fixos dos dados
originais, pois criam as próprias linhas e as removem ao final. O \cod{README.md} do pacote detalha os
pré-requisitos e os comandos.

\subsection{Princípios adotados}

\begin{itemize}
  \item \textbf{Generalidade.} As rotinas da Parte~1 recebem schema, tabela e coluna como parâmetros e obtêm tudo
        do catálogo; nenhuma delas contém nomes de tabelas, colunas ou índices do Chinook.
  \item \textbf{SQL dinâmico seguro.} Identificadores são montados com \cod{format('\%I', \dots)} e literais com
        \cod{\%L}, nunca por concatenação, o que evita injeção de SQL e erros com nomes que exigem aspas.
  \item \textbf{Verificação automática.} Cada questão possui casos de teste positivos e, onde a questão admite
        rejeição, negativos.
\end{itemize}

% =====================================================================
\section{Parte 1}

\subsection{Questão 1 --- Índices pelo catálogo}

\begin{enunciado}
Consultar as tabelas de catálogo para listar todos os índices existentes acompanhados das tabelas e colunas
indexadas pelo mesmo.
\end{enunciado}

A função \cod{bd2\_meta.listar\_indices(p\_schema, p\_incluir\_sistema)} consulta diretamente o catálogo do
PostgreSQL. A tabela \cod{pg\_index} contém uma linha por índice, com a tabela indexada (\cod{indrelid}), as
posições das colunas e os indicadores de unicidade, chave primária, exclusão e predicado parcial. Ela é combinada
com \cod{pg\_class} (para o índice e para a tabela), \cod{pg\_namespace} (schema) e \cod{pg\_am} (método de
acesso):

\begin{lstlisting}
FROM pg_index i
JOIN pg_class     idx ON idx.oid = i.indexrelid
JOIN pg_class     tbl ON tbl.oid = i.indrelid
JOIN pg_namespace ns  ON ns.oid  = tbl.relnamespace
JOIN pg_am        am  ON am.oid  = idx.relam
\end{lstlisting}

As colunas de cada posição do índice são obtidas com \cod{pg\_get\_indexdef(indexrelid, k, true)}, que devolve o
nome da coluna ou, em índices sobre expressões, a própria expressão. As colunas-chave (posições
\cod{1..indnkeyatts}) são separadas das colunas incluídas com \cod{INCLUDE}. Sem parâmetros, a função lista os
índices de tabelas e visões materializadas de todos os schemas de usuário; com \cod{p\_incluir\_sistema = true},
lista também os do catálogo do sistema e das tabelas TOAST. A visão \cod{bd2\_meta.vw\_indices} expõe o mesmo resultado.

\paragraph{Resultado.} No schema \cod{public} do Chinook a função lista \numIndicesPublic{} índices, com tabela,
colunas, método de acesso e tipo (chave primária ou comum). Os testes verificam a exclusão dos schemas de sistema
e, num índice de teste composto, parcial e com \cod{INCLUDE}, a separação correta entre colunas-chave e incluídas
(\cod{sql/06\_testes/01\_testes\_parte1\_q1\_a\_q4.sql}). A listagem completa está em
\cod{evidencias/p1\_q1\_indices.txt}.

\subsection{Questão 2 --- Procedimento que remove os índices de uma tabela}

\begin{enunciado}
Criar usando a linguagem de programação do SGBD escolhido um procedimento que remova todos os índices de uma
tabela informada como parâmetro.
\end{enunciado}

O procedimento \cod{bd2\_meta.remover\_indices\_tabela(p\_schema, p\_tabela, p\_incluir\_constraints, p\_cascade)}
localiza a tabela pelo catálogo e percorre todos os seus índices. No PostgreSQL há dois tipos de índice que
exigem comandos diferentes. Os índices independentes, criados com \cod{CREATE INDEX}, são removidos com
\cod{DROP INDEX}. Os que sustentam uma restrição \cod{PRIMARY KEY}, \cod{UNIQUE} ou \cod{EXCLUDE} não podem ser
removidos dessa forma: é necessário remover a restrição com \cod{ALTER TABLE \dots{} DROP CONSTRAINT}.

Para distinguir os casos, o procedimento associa cada índice a \cod{pg\_constraint} restringindo a tabela e o tipo
da restrição:

\begin{lstlisting}
LEFT JOIN pg_constraint con
       ON con.conindid = i.indexrelid
      AND con.conrelid = i.indrelid
      AND con.contype IN ('p', 'u', 'x')
\end{lstlisting}

O filtro por \cod{conrelid} é necessário porque as chaves estrangeiras de \emph{outras} tabelas também registram
em \cod{conindid} o índice da chave primária que referenciam; sem ele, o procedimento confundiria essas restrições
com as da própria tabela.

A chamada que remove \textbf{todos} os índices da tabela, inclusive os que sustentam chave primária e
\cod{UNIQUE}, é:

\begin{lstlisting}
CALL bd2_meta.remover_indices_tabela('schema', 'tabela',
                                     p_incluir_constraints => true);
\end{lstlisting}

Por segurança, o valor padrão de \cod{p\_incluir\_constraints} é falso: sem ele, o procedimento remove apenas os
índices independentes, pois remover uma chave primária altera as regras de integridade da tabela, e não apenas
seu desempenho. Pelo mesmo motivo, \cod{CASCADE} nunca é aplicado implicitamente. Se outras tabelas possuem chave
estrangeira para a chave primária, a remoção é recusada, a menos que se informe também \cod{p\_cascade => true},
caso em que essas chaves estrangeiras são removidas junto.

\paragraph{Resultado.} A demonstração (\cod{evidencias/p1\_q2\_remocao\_indices.txt}) é feita sobre tabelas reais
do Chinook, dentro de transações desfeitas ao final. Em \cod{invoice\_line}, a chamada padrão remove os dois índices
independentes e preserva a chave primária; com \cod{p\_incluir\_constraints} a tabela fica sem nenhum índice. Em
\cod{employee}, cuja chave primária é referenciada por outras tabelas, a remoção sem \cod{p\_cascade} é recusada, e
com \cod{p\_cascade} todos os índices e as chaves estrangeiras dependentes são removidos.

\subsection{Questão 3 --- Chaves estrangeiras pelo catálogo}

\begin{enunciado}
Consultar as tabelas de catálogo para listar todas as chaves estrangeiras existentes informando as tabelas e
colunas envolvidas.
\end{enunciado}

A função \cod{bd2\_meta.listar\_foreign\_keys(p\_schema)} lê \cod{pg\_constraint} com \cod{contype = 'f'}. Cada
chave estrangeira guarda as colunas locais em \cod{conkey} e as referenciadas em \cod{confkey}, como vetores de
números de coluna. Os dois vetores são percorridos em paralelo, o que pareia corretamente as colunas pela posição
mesmo em chaves compostas, e os números são convertidos em nomes por meio de \cod{pg\_attribute}:

\begin{lstlisting}
CROSS JOIN LATERAL unnest(c.conkey, c.confkey) WITH ORDINALITY
           AS u(attnum_local, attnum_ref, pos)
JOIN pg_attribute att_local ON att_local.attrelid = c.conrelid
                           AND att_local.attnum   = u.attnum_local
\end{lstlisting}

A função informa ainda as ações \cod{ON UPDATE} e \cod{ON DELETE}, o tipo de \cod{MATCH} e se a restrição é
postergável; a visão \cod{bd2\_meta.vw\_foreign\_keys} lista as chaves de todos os schemas.

\paragraph{Resultado.} São listadas as \numFKsPublic{} chaves estrangeiras do Chinook, entre elas a
autorreferência \cod{employee.reports\_to} $\rightarrow$ \cod{employee.employee\_id}. Os testes verificam também o
pareamento correto das colunas numa chave composta. A listagem completa está em
\cod{evidencias/p1\_q3\_foreign\_keys.txt}.

\subsection{Questão 4 --- Geração dinâmica de \texttt{CREATE TABLE}}

\begin{enunciado}
Criar usando a linguagem de programação do SGBD escolhido um script que construa de forma dinâmica a partir do
catálogo os comandos create table das tabelas existentes no esquema exemplo considerando pelo menos as
informações sobre colunas (nome, tipo e obrigatoriedade) e chaves primárias e estrangeiras.
\end{enunciado}

A função \cod{bd2\_meta.gerar\_ddl\_tabelas(p\_schema\_origem, p\_schema\_destino, p\_incluir\_fks)} devolve como
texto o script de criação de todas as tabelas de um schema; ela não executa os comandos. O schema de destino é um
parâmetro, de modo que o script pode recriar a estrutura em outro local. A geração ocorre em duas fases:

\begin{enumerate}
  \item \textbf{Tabelas e chaves primárias.} Para cada tabela, as colunas são lidas de \cod{pg\_attribute} na
        ordem de definição: nome, tipo exato com tamanho e precisão, obtido com
        \cod{format\_type(atttypid, atttypmod)}, e obrigatoriedade (\cod{attnotnull}). A chave primária vem de
        \cod{pg\_constraint} e é declarada dentro do \cod{CREATE TABLE}.
  \item \textbf{Chaves estrangeiras.} Depois de todas as tabelas, cada chave estrangeira é emitida como
        \cod{ALTER TABLE \dots{} ADD CONSTRAINT \dots{} FOREIGN KEY}, com as colunas montadas a partir de
        \cod{conkey} e \cod{confkey} e com as ações referenciais preservadas. Referências ao próprio schema de
        origem são redirecionadas ao schema de destino.
\end{enumerate}

Gerar as chaves estrangeiras numa segunda fase, como faz o utilitário \cod{pg\_dump}, dispensa ordenar as tabelas
por dependência e funciona também com autorreferências e ciclos entre tabelas. Um trecho do resultado:

\begin{lstlisting}
CREATE TABLE public.album (
    album_id integer NOT NULL,
    title character varying(160) NOT NULL,
    artist_id integer NOT NULL,
    CONSTRAINT album_pkey PRIMARY KEY (album_id)
);
...
ALTER TABLE public.album ADD CONSTRAINT album_artist_id_fkey
  FOREIGN KEY (artist_id) REFERENCES public.artist (artist_id);
\end{lstlisting}

\paragraph{Resultado.} O script completo gerado para o Chinook está em \cod{evidencias/p1\_q4\_ddl\_gerado.sql}.
Como ele reflete o catálogo atual, inclui a coluna \cod{invoice.status} criada na Questão~5. Para comprovar que o
texto gerado é válido, ele foi executado num schema novo, e a estrutura resultante coincide com a original:
\numQquatroTabelas{} tabelas, \numQquatroPKs{} chaves primárias, \numQquatroFKs{} chaves estrangeiras,
\numQquatroColunas{} colunas, das quais \numQquatroNotNull{} obrigatórias. O teste correspondente verifica ainda
que a chave estrangeira copiada rejeita uma linha órfã.

\subsection{Questão 5 --- Validação genérica de diagrama de transição de estados}

\begin{enunciado}
Implemente uma solução através da programação em banco de dados para validar os valores de uma coluna que
represente uma situação (estado) garantindo que os seus valores e suas transições atendam a especificação de um
diagrama de transição de estados (DTE). Quanto mais genérica e reutilizável for a solução melhor a pontuação
nessa questão. Junto da solução deverá ser entregue um cenário de teste demonstrando o funcionamento da solução.
\end{enunciado}

\paragraph{Ideia central.} Em vez de um trigger específico para cada tabela, o grupo implementou um motor
genérico em que \emph{o diagrama é dado, e não código}. Os estados e as transições de cada diagrama são
cadastrados em tabelas de configuração, e uma única função de trigger valida qualquer coluna de estado de
qualquer tabela consultando essa configuração. Colocar uma nova coluna sob controle de um diagrama exige apenas
cadastrá-lo e chamar um procedimento; nenhum código é escrito ou alterado.

\paragraph{Modelo.} O schema \cod{bd2\_meta} contém quatro tabelas:
\begin{itemize}
  \item \cod{state\_machine}: um diagrama e sua política de inserção. Com \cod{INITIAL\_ONLY}, uma linha nova
        precisa nascer no estado inicial; com \cod{ANY\_VALID\_STATE}, pode nascer em qualquer estado válido,
        o que é útil em cargas de dados;
  \item \cod{state\_definition}: os estados, com indicação de inicial, terminal e ativo. Um índice único parcial
        garante no máximo um estado inicial por diagrama;
  \item \cod{state\_transition}: as transições permitidas. Um trigger impede cadastrar transição que parta de
        estado terminal;
  \item \cod{state\_binding}: o vínculo entre um diagrama e uma coluna (schema, tabela e coluna).
\end{itemize}

\paragraph{Validação.} O PostgreSQL não possui trigger global; por isso a \emph{função} é única e a
\emph{instalação} é feita por tabela. A função \cod{bd2\_meta.fn\_validar\_transicao\_estado()} identifica em
tempo de execução o que está validando: a tabela vem das variáveis \cod{TG\_TABLE\_SCHEMA} e \cod{TG\_TABLE\_NAME},
a coluna vem do argumento do trigger e o valor é lido de forma genérica com \cod{to\_jsonb(NEW) -{}>{}> coluna}.
As regras são:
\begin{itemize}
  \item \textbf{inserção:} o estado não pode ser nulo, deve existir e estar ativo; com \cod{INITIAL\_ONLY}, deve
        ser o estado inicial;
  \item \textbf{atualização:} se o estado não mudou, a operação é aceita; caso contrário, o estado atual e o novo
        devem estar ativos, o estado atual não pode ser terminal e a transição (atual, novo) deve estar cadastrada;
  \item \textbf{falha segura:} se o vínculo ou o diagrama estiverem inativos, a operação é rejeitada em vez de
        prosseguir sem validação.
\end{itemize}

O procedimento \cod{bd2\_meta.registrar\_binding} verifica no catálogo a tabela, a coluna e o tipo, confere se os
dados já existentes estão em estados válidos e instala o trigger. Ela é idempotente: chamá-la novamente para a
mesma coluna atualiza o vínculo sem duplicá-lo. O procedimento \cod{remover\_binding} desfaz a instalação. O uso
completo para uma coluna nova resume-se a:

\begin{lstlisting}
INSERT INTO bd2_meta.state_machine (machine_name, insert_policy)
     VALUES ('INVOICE_STATE', 'INITIAL_ONLY');
-- INSERTs em state_definition e state_transition (estados e arestas)
CALL bd2_meta.registrar_binding('public', 'invoice', 'status',
                                'INVOICE_STATE');
\end{lstlisting}

\paragraph{Cenário de demonstração.} Para comprovar a generalidade, a mesma função, sem nenhuma alteração, é
aplicada a duas tabelas, com colunas de nomes diferentes e diagramas diferentes
(Figura~\ref{fig:dte}): \cod{public.invoice.status}, com política \cod{INITIAL\_ONLY}, e
\cod{bd2\_demo.pedido\_demo.situacao}, com política \cod{ANY\_VALID\_STATE}.

\begin{figure}[!htb]
\centering\small
\begin{tikzpicture}[->,>={Stealth},
                    every state/.style={rectangle,rounded corners,minimum height=2.2em,minimum width=6.5em,draw}]
  \node[state,initial,initial text={}] (P) at (0,0) {PENDENTE};
  \node[state,accepting] (G) at (5,0.7)  {PAGO};
  \node[state,accepting] (C) at (5,-0.7) {CANCELADO};
  \path (P.east) edge (G.west) (P.east) edge (C.west);
\end{tikzpicture}\\[0.4em]
(a) \texttt{INVOICE\_STATE}, em \texttt{public.invoice.status}\\[1.6em]
\begin{tikzpicture}[->,>={Stealth},
                    every state/.style={rectangle,rounded corners,minimum height=2.2em,minimum width=5.5em,draw}]
  \node[state,initial,initial text={}] (N) at (0,0) {NOVO};
  \node[state] (E) at (3.6,0) {EM\_ANDAMENTO};
  \node[state,accepting] (F) at (7.4,0) {CONCLUIDO};
  \node[state,accepting] (X) at (3.6,-1.6) {CANCELADO};
  \path (N) edge (E) (E) edge (F) (N) edge (X) (E) edge (X);
\end{tikzpicture}\\[0.4em]
(b) \texttt{PEDIDO\_STATE}, em \texttt{bd2\_demo.pedido\_demo.situacao}
\caption{Diagramas da demonstração. A seta de entrada indica o estado inicial; a borda dupla, os estados
terminais.}
\label{fig:dte}
\end{figure}

\paragraph{Cenário de teste.} O cenário (\cod{sql/06\_testes/02\_testes\_maquina\_estados.sql}) executa
\numCasosDTE{} casos, todos com o resultado esperado. Cada caso registra a operação, o resultado esperado e o
obtido e, nas rejeições, o código de erro. A Tabela~\ref{tab:dte} resume os grupos; a lista completa está em
\cod{evidencias/p1\_q5\_cenario\_teste.txt}.

\begin{table}[!htb]
\centering\small
\begin{tabularx}{\textwidth}{@{}lXr@{}}
\toprule
Grupo & O que verifica & Casos \\
\midrule
A & Faturas: nascimento no estado inicial, transições permitidas e proibidas, saída de estado terminal,
    estado inexistente, atualização de outras colunas & \numDTEa \\
B & Pedidos (outra tabela, outra coluna, outra política): estado nulo, estado de outro diagrama, salto de etapa,
    atualização de várias linhas desfeita por completo quando uma transição é inválida & \numDTEb \\
C & Estado desativado; vínculo ou diagrama inativo (falha segura) & \numDTEc \\
D & Registro e remoção de vínculos: idempotência, tabela, coluna ou tipo inválidos, recusa de dados existentes
    fora do diagrama & \numDTEd \\
E & Integridade do próprio diagrama: transição a partir de estado terminal, dois estados iniciais, estado
    inexistente & \numDTEe \\
F & Generalidade: as duas tabelas usam a mesma função, cujo código não contém nomes de tabela, coluna ou estado
    & \numDTEf \\
\bottomrule
\end{tabularx}
\caption{Grupos de casos do cenário de teste da Questão~5.}
\label{tab:dte}
\end{table}

% =====================================================================
\section{Parte 2}

\subsection{Questão 1 --- Regras semânticas}

\begin{enunciado}
Criar regras semânticas, que são regras que não podem ser garantidas pela estrutura do modelo relacional, usando
o esquema exemplo fornecido. As regras criadas também devem ser descritas textualmente no trabalho a ser
entregue.
\end{enunciado}

O grupo definiu três regras. Nenhuma pode ser expressa pelos recursos declarativos do modelo relacional (chave
primária, chave estrangeira, \cod{NOT NULL}, \cod{UNIQUE}, \cod{CHECK}), porque todas dependem de outras linhas
ou de outras tabelas, enquanto uma restrição \cod{CHECK} examina apenas a própria linha.

\paragraph{R-TRG-01 --- Hierarquia de funcionários sem ciclos.} Na tabela \cod{employee}, a coluna
\cod{reports\_to} indica o superior de cada funcionário. A hierarquia formada deve ser acíclica: nenhum
funcionário pode ser, direta ou indiretamente, superior de si mesmo. A chave estrangeira garante apenas que o
superior existe, e uma restrição \cod{CHECK} impediria somente a autorreferência direta; detectar um ciclo
indireto, como $A \rightarrow B \rightarrow C \rightarrow A$, exige percorrer outras linhas da tabela.

\paragraph{R-SP-01 --- Fatura nasce com pelo menos um item.} Toda fatura criada pelo usuário da aplicação deve
conter, já no momento da criação, pelo menos um item com faixa existente e quantidade inteira positiva, cobrado
pelo preço oficial da faixa. A chave estrangeira está do lado dos itens (\cod{invoice\_line} $\rightarrow$
\cod{invoice}): impede item sem fatura, mas não fatura sem itens. Uma cardinalidade mínima de um no lado
``muitos'' não é expressável declarativamente.

\paragraph{R-TRG-02 --- Total da fatura igual à soma dos itens.} Para toda fatura, \cod{invoice.total} deve ser
igual à soma de \cod{unit\_price} $\times$ \cod{quantity} dos seus itens, ou zero se ela não tiver itens. É a regra
proposta na Questão~4. Restrições \cod{CHECK} e colunas geradas não podem consultar outra tabela nem agregar
linhas.

As regras R-TRG-01 e R-TRG-02 são garantidas por triggers e a regra R-SP-01 por procedimento armazenado, atendendo ao
mínimo de um caso de cada tipo.

\subsection{Questão 2 --- Validação por trigger}

\begin{enunciado}
Implementar triggers que garantam a validação das regras semânticas criadas.
\end{enunciado}

A regra R-TRG-01 é garantida pelo trigger \cod{trg\_employee\_sem\_ciclo}, executado antes de cada inserção e de
cada alteração de \cod{reports\_to} (ou do próprio identificador do funcionário). A função aceita a operação se o funcionário passa a não ter superior ou se a
hierarquia não mudou e rejeita a autorreferência. Nos demais casos, percorre a cadeia de superiores a partir do
novo superior com uma consulta recursiva e rejeita a operação se encontrar o próprio funcionário:

\begin{lstlisting}
WITH RECURSIVE cadeia AS (
    SELECT e.employee_id, e.reports_to, ARRAY[e.employee_id] AS caminho
      FROM public.employee e
     WHERE e.employee_id = NEW.reports_to
    UNION ALL
    SELECT e.employee_id, e.reports_to, c.caminho || e.employee_id
      FROM public.employee e JOIN cadeia c ON e.employee_id = c.reports_to
     WHERE e.employee_id <> ALL (c.caminho)
)
SELECT array_agg(employee_id),
       bool_or(employee_id = NEW.employee_id OR reports_to = NEW.employee_id)
  INTO v_cadeia, v_ciclo
  FROM cadeia;
\end{lstlisting}

A segunda condição (algum elo \emph{se reporta} ao funcionário) detecta também ciclos formados por várias linhas
inseridas ou alteradas num mesmo comando.

\paragraph{Concorrência.} Duas transações simultâneas, uma tornando $A$ subordinado a $B$ e outra $B$ subordinado
a $A$, poderiam cada uma verificar a hierarquia antiga, não encontrar ciclo e confirmar. Para evitar isso, a
função obtém um bloqueio consultivo de transação (\cod{pg\_advisory\_xact\_lock}) antes da verificação, o que
serializa as alterações de hierarquia. No nível de isolamento padrão (\cod{READ COMMITTED}), a consulta seguinte
ao bloqueio enxerga a alteração já confirmada pela outra transação. Nos níveis \cod{REPEATABLE READ} e
\cod{SERIALIZABLE}, em que a fotografia dos dados é fixa, a função também bloqueia a cadeia lida com
\cod{FOR SHARE}; se um desses superiores foi alterado por outra transação, o PostgreSQL aborta a transação com
erro de serialização, que a aplicação deve repetir.

A regra R-TRG-02, também garantida por triggers, é descrita na Questão~4.

\paragraph{Resultado.} Os \numCasosH{} casos de teste (\cod{evidencias/p2\_testes\_regras.txt}) incluem autorreferência, ciclos de duas e de três pessoas,
ciclos criados por uma única instrução de várias linhas, alterações válidas e uma tentativa de ciclo na hierarquia
original do Chinook. Os testes de concorrência C3 e C4 (Seção~\ref{sec:conc}) mostram, nos dois níveis de
isolamento, a segunda sessão sendo rejeitada e a hierarquia final sem ciclo.

\subsection{Questão 3 --- Validação por procedimento armazenado e permissões}

\begin{enunciado}
Implementar procedimentos armazenados (stored procedures) que garantam a validação das regras semânticas criadas.
Nesse caso, o mecanismo de permissões deve ser utilizado para criar um usuário que somente tenha acesso à
manipulação dos dados envolvidos através do procedimento definido.
\end{enunciado}

\paragraph{Procedimento.} A regra R-SP-01 é garantida pelo procedimento \cod{public.criar\_fatura\_com\_itens}, que
recebe o cliente e a lista de itens em JSON e devolve o número da fatura criada:

\begin{lstlisting}
CALL public.criar_fatura_com_itens(1,
     '[{"track_id": 1, "quantity": 2}, {"track_id": 5, "quantity": 1}]');
\end{lstlisting}

Antes de gravar qualquer dado, o procedimento exige uma lista não vazia e valida todos os itens (identificador da
faixa e quantidade inteiros, quantidade positiva, faixa existente). Em seguida grava o cabeçalho, com os dados de
cobrança copiados do cadastro do cliente, e os itens, com o preço oficial da faixa e não um preço informado pelo
chamador. Por fim confere que todos os itens foram gravados. Não há confirmação parcial: qualquer erro desfaz o
cabeçalho e os itens. O total é calculado pelos triggers da regra R-TRG-02.

\paragraph{Usuário restrito.} Foram criados dois papéis (Tabela~\ref{tab:papeis}). O procedimento é declarado
\cod{SECURITY DEFINER} e pertence a \cod{bd2\_owner}: quando \cod{bd2\_app} o executa, ele opera com os
privilégios de \cod{bd2\_owner}, e não com os de quem o chamou. Assim, o usuário da aplicação consegue criar
faturas somente pelo caminho que valida a regra. Seguindo a recomendação da documentação do PostgreSQL para
funções \cod{SECURITY DEFINER}, o caminho de busca de objetos é fixado
(\cod{search\_path = pg\_catalog, public, pg\_temp}) e todos os objetos são qualificados pelo schema, o que
impede o chamador de substituí-los por objetos próprios. O direito de execução é retirado de \cod{PUBLIC} e
concedido apenas a \cod{bd2\_app}.

\begin{table}[!htb]
\centering\small
\begin{tabularx}{\textwidth}{@{}llX@{}}
\toprule
Papel & Atributos & Privilégios \\
\midrule
\cod{bd2\_app} & com login & Execução do procedimento. Nenhum privilégio de leitura ou escrita sobre
\cod{invoice}, \cod{invoice\_line} ou qualquer outra tabela do Chinook. \\
\cod{bd2\_owner} & sem login & Dono do procedimento. Leitura de \cod{customer} e \cod{track}; leitura e inserção
em \cod{invoice} e \cod{invoice\_line}; atualização apenas de \cod{invoice.total} (usada pela regra
R-TRG-02); uso das sequências das chaves primárias. Não pode excluir dados. \\
\bottomrule
\end{tabularx}
\caption{Papéis criados para a regra R-SP-01.}
\label{tab:papeis}
\end{table}

O usuário \cod{bd2\_app} recebe uma senha de demonstração (\cod{bd2\_app}), o que permite conectar-se de fato
como esse usuário.

\paragraph{Resultado.} Os \numCasosS{} testes executados como \cod{bd2\_app} (\cod{evidencias/p2\_testes\_regras.txt})
confirmam que o procedimento cria a fatura com os
itens e o total corretos e rejeita lista vazia ou malformada, cliente ou faixa inexistentes e quantidades iguais a zero,
negativas ou fracionárias, sem deixar nenhuma fatura parcial gravada. Inserção, alteração, exclusão,
\cod{TRUNCATE} e mesmo leitura diretos nas tabelas de faturas são negados por falta de permissão. Outros \numCasosP{} casos
conferem no catálogo os atributos dos papéis e do procedimento (\cod{evidencias/p2\_permissoes.txt} mostra os
privilégios efetivos). No teste L1 (Seção~\ref{sec:conc}), uma conexão
real como \cod{bd2\_app}, autenticada por senha, cria uma fatura pelo procedimento e tem a inserção e a leitura
diretas negadas; uma senha errada é recusada.

\subsection{Questão 4 --- Integridade de \texttt{Invoice.Total}}

\begin{enunciado}
A base original do Chinook possui uma coluna Total na tabela Invoice representada de forma redundante com as
informações contidas nas colunas UnitPrice e Quantity na tabela InvoiceLine. Podemos identificar nesse caso uma
regra semântica onde o valor Total de um Invoice deve ser igual à soma de UnitPrice * Quantity de todos os
registros de InvoiceLine relacionados a um Invoice. Implementar uma solução que garanta a integridade dessa
regra.
\end{enunciado}

A solução faz com que o total passe a ser mantido exclusivamente pelo banco. Antes de ativar a regra, o script
confere que os dados originais a respeitam; no Chinook~1.4.5 não há divergências. A regra é garantida por três
triggers:

\begin{table}[!htb]
\centering\small
\begin{tabularx}{\textwidth}{@{}lp{4.6cm}X@{}}
\toprule
Trigger & Quando dispara & O que faz \\
\midrule
A & antes de inserir uma fatura ou alterar seu total (ou seu identificador) & Substitui o valor informado pela soma calculada dos
itens. Um total digitado manualmente é corrigido. \\
B & após inserir, excluir ou alterar fatura, preço ou quantidade de um item & Recalcula a fatura do item e, se o
item mudou de fatura, também a anterior, sempre em ordem crescente de identificador para evitar impasse entre
transações. \\
C & após \cod{TRUNCATE} dos itens & Zera o total de todas as faturas, já que esse comando não dispara triggers
por linha. \\
\bottomrule
\end{tabularx}
\caption{Triggers da regra R-TRG-02.}
\end{table}

Com isso ficam cobertos a inclusão e a exclusão de itens, a alteração de preço, quantidade ou fatura de um item,
a inclusão de fatura, a alteração direta do total, o esvaziamento da tabela de itens e a instrução
\cod{INSERT \dots{} ON CONFLICT DO UPDATE}.

\paragraph{Resultado.} Os \numCasosT{} casos de teste (\cod{evidencias/p2\_testes\_regras.txt}) conferem o total após cada um desses eventos, inclusive a
transferência de um item entre duas faturas (ambas recalculadas) e a inserção de cinquenta itens numa única
instrução. Ao final, uma auditoria sobre todas as faturas não encontra nenhuma divergência.

\subsection{Testes de concorrência}
\label{sec:conc}

Um teste executado numa única sessão não demonstra o comportamento concorrente. Por isso o script
\cod{sql/06\_testes/04\_concorrencia.sh} abre conexões \cod{psql} reais em paralelo
(\cod{evidencias/p2\_concorrencia.txt}):

\begin{itemize}
  \item \textbf{C1:} \numSessoes{} sessões inserem itens na mesma fatura ao mesmo tempo; o total final é igual à
        soma, sem erros;
  \item \textbf{C2:} \numSessoes{} sessões transferem itens entre duas faturas em sentidos opostos; não ocorre
        impasse e os dois totais terminam corretos;
  \item \textbf{C3 e C4:} duas sessões tentam fechar um ciclo na hierarquia, em \cod{READ COMMITTED} e em
        \cod{REPEATABLE READ}; nos dois casos a segunda é rejeitada e a hierarquia final não tem ciclo;
  \item \textbf{C5:} \numSessoes{} sessões, todas com o papel \cod{bd2\_app} (\cod{SET ROLE}), executam o procedimento
        simultaneamente e todas as faturas são criadas corretamente;
  \item \textbf{L1:} conexão real como \cod{bd2\_app}, descrita na Questão~3.
\end{itemize}

% =====================================================================
\section{Limitações e decisões de projeto}

\begin{itemize}
  \item \textbf{Remoção de índices.} A remoção de índices que sustentam restrições e o uso de \cod{CASCADE} são
        opcionais e precisam ser pedidos explicitamente, para evitar a perda acidental de regras de integridade.
  \item \textbf{Escopo do gerador de DDL.} O gerador cobre o exigido: colunas com nome, tipo e obrigatoriedade,
        chaves primárias e estrangeiras. Não gera valores padrão (como as sequências da variante utilizada),
        restrições \cod{UNIQUE} e \cod{CHECK}, índices nem comentários. Por isso, um esquema em que uma chave
        estrangeira referencie uma restrição \cod{UNIQUE}, e não a chave primária, não seria reproduzido; isso não
        ocorre no Chinook.
  \item \textbf{Tipos de coluna de estado.} O motor aceita colunas de texto, enumeração e inteiras. Colunas
        \cod{char(n)} são recusadas, porque o preenchimento com espaços prejudica a comparação dos estados. Os vínculos são
        guardados por nome: se a tabela ou a coluna vinculada for renomeada, o motor passa a rejeitar as operações
        nessa tabela (falha segura) até que o vínculo seja refeito.
  \item \textbf{Exclusão de linhas.} Excluir uma linha não é uma transição de estado e não é validado pelo motor.
        Num diagrama com política \cod{ANY\_VALID\_STATE}, excluir e reinserir equivale a recriar a linha em outro
        estado; quando isso não for desejado, deve-se usar \cod{INITIAL\_ONLY} ou retirar o direito de exclusão.
  \item \textbf{Total corrigido, e não rejeitado.} Na regra R-TRG-02, um total informado manualmente é substituído
        pelo valor calculado em vez de provocar erro, o que mantém o total sempre correto sem exigir que as
        aplicações o calculem.
  \item \textbf{Alcance das regras.} A regra R-SP-01 vale para o usuário da aplicação, cujo único caminho de
        escrita é o procedimento; usuários administrativos, com acesso direto às tabelas, não passam por ele.
        Superusuários e donos das tabelas podem ainda desabilitar triggers.
  \item \textbf{Senha de demonstração.} Por se tratar de um trabalho acadêmico, a senha de \cod{bd2\_app} é fixa
        e pode ser alterada pela variável \cod{BD2\_APP\_PASSWORD}; em produção, ela não ficaria no código.
\end{itemize}

% =====================================================================
\section{Conclusão}

Todas as questões das duas partes foram implementadas em PL/pgSQL sobre o Chinook. As rotinas de catálogo e o
motor de estados da Parte~1 são genéricos e não dependem do esquema exemplo; o motor, em particular, permite
controlar qualquer coluna de estado apenas com dados de configuração. Na Parte~2, as três regras semânticas foram
garantidas por triggers e por um procedimento armazenado combinado a permissões, considerando inclusive o acesso
concorrente. A execução completa, a partir de um banco limpo, aplica todos os scripts e confirma o comportamento
descrito por meio de testes automatizados, cujas saídas acompanham a entrega.

\end{document}
